\documentclass[twoside]{article}

\usepackage[preprint]{aistats2027}
\usepackage[round,authoryear]{natbib}
\usepackage{amsmath,amssymb,amsthm}
\usepackage{mathtools}
\usepackage{microtype}
\usepackage{booktabs}
\usepackage{url}

\renewcommand{\bibname}{References}
\renewcommand{\bibsection}{\subsubsection*{\bibname}}
\bibliographystyle{apalike}

\newtheorem{theorem}{Theorem}
\newtheorem{proposition}[theorem]{Proposition}
\newtheorem{corollary}[theorem]{Corollary}
\newtheorem{lemma}[theorem]{Lemma}
\newtheorem{remark}[theorem]{Remark}
\newcommand{\E}{\mathbb E}
\newcommand{\PP}{\mathbb P}
\newcommand{\cF}{\mathcal F}
\newcommand{\cX}{\mathcal X}
\newcommand{\TV}{\operatorname{TV}}
\newcommand{\err}{\operatorname{err}}
\newcommand{\Ber}{\operatorname{Ber}}
\newcommand{\Poi}{\operatorname{Poi}}
\newcommand{\1}{\mathbf 1}

\begin{document}
\runningauthor{}

\twocolumn[
\aistatstitle{When Sharing Fails Under Massart Noise: Near-Linear Lower Bounds for Multi-Distribution Learning}

\aistatsauthor{Pahan Dewasurendra}

\aistatsaddress{Johns Hopkins University}
]

\begin{abstract}
Personalized multi-distribution learning asks for one predictor per data source while exploiting a Bayes classifier shared by all sources. Under Massart noise, separate empirical risk minimization uses roughly $kd/\epsilon$ samples for $k$ sources and VC dimension $d$. The best known dimension-dependent sharing term is only $k\sqrt d/\epsilon$. Its authors ask whether higher-order moment matching can close the gap.

We give a near-linear answer in the high-accuracy regime. For noise bound $\eta<1/2$, let $\gamma=1-2\eta$ and $z_\eta=2(1-\eta)/\eta-1$. Define \[
 \begin{aligned}
 K_m(z)&=1+2\sum_{j=1}^m T_j(z)^2,\\
 w_m&=K_m(z_\eta)^{-1},
 \end{aligned}
\] where $T_j$ is a Chebyshev polynomial. Subject to explicit feasibility and concentration conditions, every adaptive learner requires \[ \Omega\!\left( \frac{k m\gamma w_m}{\epsilon} d^{\frac{2m}{2m+1}} \right) \] samples. Consequently, for every fixed $\zeta>0$ the lower bound is $\Omega_{\eta,\zeta}(k d^{1-\zeta}/\epsilon)$ at sufficiently high accuracy. If $\epsilon$ is polynomially small in $d$, it becomes $\Omega_\eta(kd\exp[-C_\eta\sqrt{\log d}]/\epsilon)$. Thus sharing cannot yield a polynomial improvement over separate learning in this regime.

The key is an exact external-atom moment theorem. An arcsine law supported on label-zero biases admits a moment-matching measure with the maximum possible atom $1/K_m(z)$ at a prescribed label-one bias. Matching $2m$ moments delays distinguishability until a $(2m+1)$-fold coordinate collision. A randomized null and a variance-sensitive holdout then lift this single-source barrier to $k$ adaptively sampled sources without increasing VC dimension.
\end{abstract}

\section{Introduction}
\label{sec:intro}

Learning systems often draw data from several populations, environments, or clients. Multi-distribution learning formalizes the benefit of choosing which source to sample while requiring guarantees across all sources \citep{blum2017collaborative,haghtalab2022ondemand}. In the realizable setting, a shared classifier lets $k$ sources reduce their total cost from separate learning, about $kd/\epsilon$, to nearly $(d+k)/\epsilon$. Analogous additive rates are possible in agnostic minimax formulations \citep{zhang2024optimal,peng2024sample}.

Bounded label noise lies between these two regimes. Suppose every source has the same Bayes classifier $f^\star$ in a class of VC dimension $d$, and its pointwise label error is at most $\eta<1/2$. A single source can be learned at the fast rate $\widetilde O(d/[\epsilon(1-2\eta)])$. It is natural to ask whether the shared Bayes rule preserves the realizable benefit across sources.

\citet{hanashiro2026bounded} recently showed that the answer depends on the benchmark. If every personalized output must compete with its source's unknown Bayes error, an additional $\Omega(k\sqrt d/\epsilon)$ term is unavoidable, even though the Bayes classifier is shared. Their lower bound draws latent label biases from two priors with the same first moment. One sample from a light coordinate has the same law under both priors, so a repeated coordinate is needed. The birthday threshold gives $\sqrt d$.

This leaves a wide gap from the separate-learning upper bound. The source paper explicitly suggests higher-order moment matching, but identifies two obstacles. A useful alternative must keep enough mass above $1/2$ to create a fixed excess-risk gap. It must also lift to many sources while preserving one shared classifier in a class of dimension $d$. A constant fraction of easily detectable high coordinates does not reach linear dependence.

We resolve these obstacles up to a subpolynomial factor when accuracy grows with dimension. The first ingredient is a sharp truncated moment result that may be useful beyond this problem. Begin with the arcsine law on a low-noise interval and prescribe an atom outside that interval. The largest atom compatible with the first $2m$ moments is the reciprocal Christoffel kernel at the external point. This atom decreases exponentially with $m$, at a rate that we characterize exactly. We compensate by increasing the total mass of the light coordinates. This keeps the excess-risk gap fixed while making the high coordinates rarer.

The second ingredient turns moment matching into passive sample hardness. After Poissonization, the label counts at a coordinate have likelihood proportional to an expected degree-$s+t$ polynomial of the latent bias. Matching $2m$ moments makes all count pairs with total at most $2m$ identical. Distinguishing requires a $(2m+1)$-fold collision, whose threshold is $d^{2m/(2m+1)}$ \citep{wu2020polynomial,canonne2026short}.

Finally, we use disjoint source supports but a single randomized null. A source-specific alternative changes one block, and its high-coordinate indicator remains a shared Bayes classifier. A fresh holdout converts an accurate personalized output into a test. Its variance is of order $\epsilon/(1-2\eta)$, so the holdout costs only $O(1/[\epsilon(1-2\eta)])$ samples. The resulting test forces the learner to pay the collision cost on every source under the same valid null.

\paragraph{Contributions.}
Our results are entirely information theoretic.
\begin{itemize}
\item We identify an exact external-atom construction for $2m$ truncated
moments of the arcsine law. The feasible atom has mass $1/K_m(z)$, and no larger atom is possible for this null prior.
\item We obtain a parameterized personalized Massart lower bound with
dimension exponent $2m/(2m+1)$. It applies to adaptive source selection and improper, randomized outputs.
\item Optimizing the moment order gives
$kd\exp[-O_\eta(\sqrt{\log d})]/\epsilon$ at polynomial accuracy. For any fixed exponent loss $\zeta$, a fixed-order consequence gives $k d^{1-\zeta}/\epsilon$.
\item The noise dependence is explicit through a Chebyshev kernel. This
extends the prior constant-noise lower bound and isolates how the high atom shrinks as the Massart margin changes.
\end{itemize}

\section{Personalized Massart Learning}
\label{sec:setup}

Let $P_1,\ldots,P_k$ be distributions on $\cX\times\{0,1\}$. For a binary predictor $f$, write $\err_i(f)=P_i(f(X)\ne Y)$. We assume there is an $f^\star\in\cF$ such that, for every $i$ and $x$ in the support of $P_i$,
\begin{equation}
 P_i(Y\ne f^\star(X)\mid X=x)\le\eta<\frac12.
 \label{eq:massart}
\end{equation}
Hence $f^\star$ is Bayes optimal on every source. Let $\eta_i^\star=\err_i(f^\star)$.

A personalized learner adaptively chooses source indices and receives an independent labeled sample from the chosen source. It returns arbitrary possibly randomized predictors $\widehat f_1,\ldots,\widehat f_k$. Success means
\begin{equation}
 \max_{i\in[k]}\{\err_i(\widehat f_i)-\eta_i^\star\}\le\epsilon
 \label{eq:objective}
\end{equation}
with probability at least $1-\delta$. This is the \textsc{MDL-Mass} objective of \citet{hanashiro2026bounded}. Our lower bound allows improper outputs and fully adaptive source choices.

Write $\gamma=1-2\eta$. Learning each source separately by VC empirical risk minimization gives the benchmark
\begin{equation}
 O\!\left(
 \frac{k[d\log(1/\epsilon)+\log(k/\delta)]}{\gamma\epsilon}
 \right).
 \label{eq:separate}
\end{equation}
The fast $1/\epsilon$ dependence follows from the Massart variance condition \citep{massart2006risk,boucheron2005classification}.

Let $T_j$ denote the degree-$j$ Chebyshev polynomial of the first kind. Define
\begin{align}
 z_\eta&=\frac{2(1-\eta)}{\eta}-1>1,\qquad
 \gamma=1-2\eta,\label{eq:zeta}\\
 K_m(z)&=1+2\sum_{j=1}^mT_j(z)^2,\qquad
 w_m=K_m(z_\eta)^{-1}.\label{eq:kernel}
\end{align}
The kernel has a useful closed form. With $A_\eta=\operatorname{arcosh}(z_\eta)$,
\begin{equation}
 K_m(z_\eta)=m+\frac12+
 \frac{\sinh((2m+1)A_\eta)}{2\sinh(A_\eta)}.
 \label{eq:kernel-closed}
\end{equation}
In terms of the Massart margin, $z_\eta=(1+3\gamma)/(1-\gamma)$ and $A_\eta=\sqrt{8\gamma}(1+O(\gamma))$ as $\gamma\downarrow0$. The exact form shows two regimes. If $mA_\eta\le1$, then $K_m=\Theta(m)$. If $mA_\eta\ge1$, its growing term is of order $e^{(2m+1)A_\eta}/\sinh(A_\eta)$. Thus the construction records how many moments can be matched before the prescribed label switch becomes too rare.

\begin{theorem}[Parameterized lower bound]
\label{thm:main}
There are universal constants $c,C>0$ with the following property. Fix $\eta\in(0,1/2)$ and an integer $m\ge1$. If
\begin{align}
 \epsilon&\le c\gamma w_m,\label{eq:cond-eps}\\
 dw_m&\ge C,\label{eq:cond-count}\\
 m\gamma^2w_m d^{\frac{2m}{2m+1}}&\ge C,\label{eq:cond-holdout}
\end{align}
then there is a binary class of VC dimension exactly $d$ for which every learner satisfying \eqref{eq:objective} with probability at least $0.9$ uses
\begin{equation}
 \Omega\!\left(
 \frac{k m\gamma w_m}{\epsilon}
 d^{\frac{2m}{2m+1}}
 \right)
 \label{eq:main-bound}
\end{equation}
samples in the worst case.
\end{theorem}

The three conditions have separate roles. Equation~\eqref{eq:cond-eps} keeps the hard distribution normalized. Equation~\eqref{eq:cond-count} concentrates the number of high coordinates. Equation~\eqref{eq:cond-holdout} makes the collision barrier dominate the final validation cost. All hold for the optimized choice below when $\eta$ is fixed and $\epsilon$ is polynomially small in $d$.

\begin{corollary}[No polynomial gain from sharing]
\label{cor:power}
Fix $\eta\in(0,1/2)$ and $\zeta>0$. There are positive constants $c_{\eta,\zeta}$ and $C_{\eta,\zeta}$ such that, whenever $\epsilon\le c_{\eta,\zeta}$ and $d\ge C_{\eta,\zeta}$, the lower bound is
\begin{equation}
 \Omega_{\eta,\zeta}\!\left(\frac{k d^{1-\zeta}}{\epsilon}\right).
 \label{eq:power}
\end{equation}
\end{corollary}

\begin{corollary}[Polynomial accuracy]
\label{cor:optimized}
Fix $\eta\in(0,1/2)$ and $\beta>0$. For all sufficiently large $d$, if $\epsilon\le d^{-\beta}$, every personalized learner can require
\begin{equation}
 \Omega_{\eta,\beta}\!\left(
 \frac{kd}{\epsilon}
 \exp[-C_\eta\sqrt{\log d}]
 \right)
 \label{eq:optimized}
\end{equation}
samples.
\end{corollary}

Corollary~\ref{cor:power} follows by fixing $m$ with $(2m+1)^{-1}\le\zeta$. For Corollary~\ref{cor:optimized}, take $2m\asymp\sqrt{\log(d)/A_\eta}$. Equation~\eqref{eq:kernel-closed} gives $K_m\le C_\eta e^{2mA_\eta}$. The losses $2mA_\eta$ and $\log(d)/(2m+1)$ balance at order $\sqrt{\log d}$. Polynomial accuracy makes \eqref{eq:cond-eps} eventually valid, and the other two left sides diverge. Comparison with \eqref{eq:separate} shows that the remaining gap is subpolynomial in $d$, apart from logarithms.

Table~\ref{tab:orders} displays the exact tradeoff at the canonical noise bound $\eta=0.49$ used by \citet{hanashiro2026bounded}. Each additional moment order raises the collision exponent but reduces the external atom, which in turn tightens the accuracy condition. These are evaluations of \eqref{eq:kernel}, not fitted or simulated values.

\begin{table}[t]
\centering
\caption{Moment order and hard-atom mass at $\eta=0.49$.}
\label{tab:orders}
\begin{tabular}{@{}rrrr@{}}
\toprule
$m$ & matched & $2m/(2m+1)$ & $w_m$ \\
\midrule
1 & 2  & 0.667 & 0.2994 \\
2 & 4  & 0.800 & 0.1443 \\
4 & 8  & 0.889 & 0.03710 \\
8 & 16 & 0.941 & 0.001768 \\
\bottomrule
\end{tabular}
\end{table}

\section{An Exact External Atom}
\label{sec:moment}

Let $\sigma$ be the arcsine probability measure on $[-1,1]$,
\begin{equation}
 d\sigma(u)=\frac{du}{\pi\sqrt{1-u^2}}.
 \label{eq:arcsine}
\end{equation}
Its orthonormal polynomials are $\phi_0=1$ and $\phi_j=\sqrt2T_j$ for $j\ge1$. Thus $K_m(z)$ in \eqref{eq:kernel} is its degree-$m$ Christoffel kernel evaluated at $z$.

\begin{lemma}[Exact external-atom moments]
\label{lem:moment}
For every $z>1$ and $m\ge1$, there is a probability measure $\nu_m$ supported on $[-1,1]$ such that
\begin{equation}
 \sigma_1=\frac1{K_m(z)}\delta_z+
 \left(1-\frac1{K_m(z)}\right)\nu_m
 \label{eq:moment-pair}
\end{equation}
matches the moments of $\sigma$ through degree $2m$. Moreover, $1/K_m(z)$ is the largest possible mass at $z$ among all matching probability measures whose remaining mass is supported on $[-1,1]$.
\end{lemma}

\begin{proof}
Set $w=1/K_m(z)$ and define, on polynomials of degree at most $2m$, \[ L(P)=\int P\,d\sigma-wP(z). \] Every $P\ge0$ on $[-1,1]$ has a Markov--Lukacs representation
\begin{equation}
 P(u)=s(u)^2+(1-u^2)t(u)^2,
 \label{eq:lukacs}
\end{equation}
where $\deg(s)\le m$ and $\deg(t)\le m-1$ \citep[Section 2.5]{wu2020polynomial}. The reproducing-kernel inequality gives
\begin{equation}
 s(z)^2\le K_m(z)\int s^2\,d\sigma.
 \label{eq:christoffel}
\end{equation}
Since $z>1$, equations \eqref{eq:lukacs} and \eqref{eq:christoffel} imply
\begin{align*}
 L(P)&=\int s^2d\sigma-ws(z)^2+\int(1-u^2)t^2d\sigma\\
 &\quad+w(z^2-1)t(z)^2\ge0.
\end{align*}
The truncated Riesz representation theorem gives a nonnegative representing measure on $[-1,1]$ for $L$. Its mass is $L(1)=1-w$. Normalizing proves \eqref{eq:moment-pair}.

For maximality, any matching measure with atom $w'$ at $z$ satisfies $w's(z)^2\le\int s^2d\sigma$ for every degree-$m$ polynomial $s$. Taking the reproducing kernel $s(u)=\sum_{j=0}^m\phi_j(z)\phi_j(u)$ gives $w'\le1/K_m(z)$.
\end{proof}

Map $u\mapsto\eta(u+1)/2$. The null law $\mu_0$ is supported on $[0,\eta]$, while the external point $z_\eta$ maps to $1-\eta$. Lemma \ref{lem:moment} produces
\begin{equation}
 \mu_1=w_m\delta_{1-\eta}+(1-w_m)\nu_m,
 \qquad \operatorname{supp}(\nu_m)\subseteq[0,\eta],
 \label{eq:bias-priors}
\end{equation}
with moments matching through order $2m$. Every null bias favors label zero. The distinguished alternative atom favors label one, and both have Massart margin exactly $\gamma$ at the closest endpoints.

\section{Collisions Hide the Bayes Rule}
\label{sec:testing}

Consider a domain $\{0,1,\ldots,d\}$. Put covariate mass $1-\alpha$ at zero and $\alpha/d$ at each light coordinate. The bias at zero is fixed below $1/2$. At each light coordinate $x$, draw $q_x$ independently from $\mu_h$ for $h\in\{0,1\}$, then keep $q_x$ fixed and draw $Y\mid X=x\sim\Ber(q_x)$ on every visit. This fixed latent bias is what makes coordinate collisions informative.

\begin{lemma}[Moment-collision testing]
\label{lem:testing}
Any test between these two random environments with both errors at most $1/5$ requires
\begin{equation}
 \Omega\!\left(
 \frac{m}{\alpha}d^{\frac{2m}{2m+1}}
 \right)
 \label{eq:test-lower}
\end{equation}
samples, provided the displayed quantity exceeds a universal constant.
\end{lemma}

\begin{proof}[Proof idea]
Poissonize the sample count with mean $N$. At one light coordinate, the counts of labels one and zero have conditional law \[ \Poi(\lambda q)\otimes\Poi(\lambda(1-q)), \qquad \lambda=N\alpha/d. \] The mixed probability of a count pair $(s,t)$ is
\begin{equation}
 \frac{e^{-\lambda}\lambda^{s+t}}{s!t!}
 \E[q^t(1-q)^s].
 \label{eq:count-likelihood}
\end{equation}
This expectation is a polynomial of degree $s+t$. The two laws agree for $s+t\le2m$. Their total count is $\Poi(\lambda)$ under both priors. Product subadditivity therefore gives
\begin{equation}
 \TV(Q_0,Q_1)\le d\,\PP(\Poi(\lambda)\ge2m+1).
 \label{eq:product-tv}
\end{equation}
For \[ N\le \frac{2m+1}{4e\alpha}d^{\frac{2m}{2m+1}}, \] the Poisson Chernoff bound makes \eqref{eq:product-tv} at most $4^{-(2m+1)}\le1/64$. A standard lower-tail coupling converts any fixed-budget test into a Poissonized test at constant overhead. Full constants are in the supplement.
\end{proof}

Set
\begin{equation}
 \alpha=\frac{16\epsilon}{\gamma w_m}.
 \label{eq:alpha}
\end{equation}
Condition \eqref{eq:cond-eps} ensures $\alpha\le1$. Under the null, zero is Bayes optimal. Under the alternative, let $H$ count coordinates at the atom $1-\eta$. Then $H\sim\operatorname{Bin}(d,w_m)$. Condition \eqref{eq:cond-count} and a Chernoff bound make
\begin{equation}
 \frac{dw_m}{2}\le H\le2dw_m
 \label{eq:high-event}
\end{equation}
hold with constant high probability. On this event, the high-coordinate indicator is Bayes, while zero has excess
\begin{equation}
 \Delta_0=\frac{\alpha\gamma H}{d}\in[8\epsilon,32\epsilon].
 \label{eq:gap}
\end{equation}
Combining \eqref{eq:test-lower} and \eqref{eq:alpha} gives the one-source cost $\Omega(m\gamma w_m d^{2m/(2m+1)}/\epsilon)$.

\section{Forcing the Cost on Every Source}
\label{sec:lifting}

Take $k$ disjoint blocks $\cX_1,\ldots,\cX_k$, each containing one heavy and $d$ light points. Let $\cF$ contain the zero function and every indicator of a subset of one block's light points. It shatters a light block. No set of $d+1$ points is shattered. A set containing a heavy point cannot label that point one, and a light-only set spanning two blocks cannot label points in both blocks one. Hence $\operatorname{VCdim}(\cF)=d$.

Each source is supported on its corresponding block with the marginal used in Section~\ref{sec:testing}. Under the common null, independently draw one fixed $\mu_0^d$ bias vector for every block. Every realization is a valid Massart instance with shared Bayes classifier zero.

For the source-$i$ alternative, replace only its bias prior by $\mu_1^d$. The Bayes classifier is the indicator of source $i$'s high coordinates. It belongs to $\cF$ and equals zero on every other source support, where it is also Bayes. Thus every realized alternative still has one shared classifier.

It remains to turn a learner into the test from Lemma~\ref{lem:testing}. Let $\widehat f_i$ be an $\epsilon$-optimal output and define on a fresh source-$i$ sample
\begin{equation}
 Z=\1\{f_0(X)\ne Y\}-\1\{\widehat f_i(X)\ne Y\},
 \label{eq:validation-variable}
\end{equation}
where $f_0$ is zero. Under the null, $\E Z\in[-\epsilon,0]$. Under an alternative satisfying \eqref{eq:high-event}, $\E Z\ge\Delta_0-\epsilon\ge7\epsilon$.

The Massart margin controls validation variance. For any predictor $f$ with excess at most $\epsilon$,
\begin{equation}
 P_i(f(X)\ne f^\star(X))\le\epsilon/\gamma.
 \label{eq:disagreement}
\end{equation}
On \eqref{eq:high-event}, zero disagrees with Bayes on mass at most $2\alpha w_m=32\epsilon/\gamma$. Therefore
\begin{equation}
 \operatorname{Var}(Z)\le
 P_i(f_0(X)\ne\widehat f_i(X))\le33\epsilon/\gamma.
 \label{eq:validation-variance}
\end{equation}
Bernstein's inequality estimates $\E Z$ to $2\epsilon$ from $O(1/(\gamma\epsilon))$ fresh samples. Thresholding at $3\epsilon$ separates the two priors.

\begin{proof}[Proof of Theorem~\ref{thm:main}]
Let $L$ be a sufficiently small constant times the one-source lower bound. Fix source $i$, and run the MDL learner with source $i$ supplied by an unknown testing environment. Simulate all other sources from independent null draws. If the learner asks for $L/2$ samples from source $i$, truncate and declare the alternative. If it finishes first, run the holdout test above.

Suppose the learner's expected source-$i$ count under the common null were $o(L)$. Markov's inequality makes truncation rare under the null. Learner success and holdout concentration then make the test accept with constant high probability. Under the alternative, truncation is already a correct decision. Without truncation, event \eqref{eq:high-event}, learner success, and holdout concentration make the test reject with constant high probability.

The test uses at most $L/2+O(1/(\gamma\epsilon))$ samples from its unknown oracle. Condition \eqref{eq:cond-holdout} makes this less than the threshold in Lemma~\ref{lem:testing}, a contradiction. Hence the expected count from every source is $\Omega(L)$ under the same common randomized null. Summing over $i$ proves \eqref{eq:main-bound}. Pointwise correctness of the learner justifies averaging over this null prior.
\end{proof}

\section{Context and Limitations}
\label{sec:discussion}

The moment-pair technique is classical. Polynomial approximation and moment duality underlie many sharp statistical lower bounds \citep{wu2019chebyshev,wu2020polynomial}. Recent exposition makes explicit that matching $r$ moments delays testing until an $(r+1)$-fold collision \citep{canonne2026short}. Our external-atom lemma identifies the exact atom compatible with the arcsine moments, then uses it to resolve the open dimension exponent in personalized Massart learning up to a subpolynomial factor.

Most MDL results optimize one hypothesis against the worst source \citep{haghtalab2022ondemand,zhang2024optimal,peng2024sample}. Personalized realizable learning admits additive dependence on $d$ and $k$ \citep{blum2017collaborative}. If different sources may have different target classifiers, separate-learning dependence can return \citep{deng2024different}. Here the target is shared. The hardness comes only from adapting to each source's unknown Bayes error. Active and limited-round MDL address complementary access restrictions \citep{zhang2025active,haghtalab2025adaptivity}.

The accuracy condition is substantive. At fixed $\epsilon$, increasing $m$ eventually makes the external atom too rare to maintain the risk gap while keeping $\alpha\le1$. Our theorem then gives a fixed exponent below one, not a uniform $d^{1-o(1)}$ statement. Determining the exact joint frontier for fixed accuracy remains open. The explicit kernel may help with that question.

Our lower bound is also information theoretic. It does not impose round, properness, or computational restrictions, and it does not propose a better algorithm. Its comparison point is separate ERM, which already matches the new lower bound up to a subpolynomial dimension factor in the stated regime.

\bibliography{references}

\clearpage
\end{document}
